%%%%%%%%%%%%%%%%%%%%%%%%%%%%%%%%%%%%%%%%%%%%%%%%%%%%%%%%%%%%%%%
%
% Welcome to Overleaf --- just edit your LaTeX on the left,
% and we'll compile it for you on the right. If you give
% someone the link to this page, they can edit at the same
% time. See the help menu above for more info. Enjoy!
%
%%%%%%%%%%%%%%%%%%%%%%%%%%%%%%%%%%%%%%%%%%%%%%%%%%%%%%%%%%%%%%%

%% BIO570 Project Proposal LaTeX template
%% Based on the EU FET Open  Proposal LaTeX template and h2020proposal.cls LaTeX class for writing EU H2020 RIA proposals.
%% 
%% Copyright (c) 2017, Tia Tidwell
%%
%%  This latex class is free software: you can redistribute it and/or modify
%%  it under the terms of the GNU General Public License as published by
%%  the Free Software Foundation, either version 3 of the License, or
%%  (at your option) any later version.
%%
%%  h2020proposal.cls is distributed in the hope that it will be useful,
%%  but WITHOUT ANY WARRANTY; without even the implied warranty of
%%  MERCHANTABILITY or FITNESS FOR A PARTICULAR PURPOSE.  See the
%%  GNU General Public License for more details.
%%
%%  You should have received a copy of the GNU General Public License
%%  along with Foobar.  If not, see <http://www.gnu.org/licenses/>.
%%
%%
%%
%% Disclaimer: The template is based on the document provided by the EU Participants Portal 
%% at the address: http://ec.europa.eu/research/participants/data/ref/h2020/call_ptef/pt/h2020-call-pt-ria-ia_en.pdf
%%
\documentclass[11pt]{article} % Loads settings for the document layout

\usepackage{geometry}
 \geometry{
 a4paper,
 total={210mm,297mm},
 left=25mm,
 right=25mm,
 top=25mm,
 bottom=25mm,
 }
\usepackage{setspace}
 \doublespacing
 
\usepackage{gensymb}

\usepackage{graphicx}

% Preamble

% The following settings are used for title generation and will show up in the
% main document where the \maketitle command is set.
\title{Interesting and descriptive proposal project title}
\date{\today}
\author{Author 1, Author 2, \& Author 3}

% Main document

\begin{document} % The document starts here

\maketitle % Creates the titlepage

\pagenumbering{arabic} % Turns on page numbering

\noindent\rule{\linewidth}{\arrayrulewidth}

\section*{Project Summary}
Limit to 250 words. The Project
Summary consists of an overview, a statement on the intellectual merit of the proposed activity, and a statement
on the impacts of the proposed activity.
The overview includes a description of the activity that would result if the proposal were funded and a statement
of objectives and methods to be employed. 

\noindent\rule{\linewidth}{\arrayrulewidth}

\clearpage
\tableofcontents
\noindent\rule{\linewidth}{\arrayrulewidth}


\section{Introduction}
\subsection{Background}
Give relevant information to put your study in context of current knowledge and introduce the problem to which you are responding.

\subsection{Aims/Objectives}
\begin{itemize}
\item Describe the targeted scientific breakthrough of the project.
\item Describe the specific objectives for the project, which should be clear, measurable, realistic and achievable within the duration of the project.
\end{itemize}

\subsection{Impact/Significance}
\begin{itemize}
\item Describe the advance your proposal would provide beyond the state-of-the-art, and to what extent the proposed work is ambitious, novel and of a foundational nature. Your answer could refer to the ground-breaking nature of the objectives, concepts involved, issues and problems to be addressed, and approaches and methods to be used. 
\item Describe the importance of the technological outcome with regards to its transformational impact on science, technology and/or society.
\end{itemize}

\section{Experimental design and methods}
Describe the overall research approach, the methods, and explain relevance to the objectives.  


\subsection{Dissemination and exploitation of results}
\begin{itemize}
\item Provide a plan for disseminating and exploiting the project results. The plan, which should be proportionate to the scale of the project, should contain measures to be implemented both during and after the project. 
\item Explain how the proposed measures will help to achieve the expected impact of the project. 
\item Where relevant, include information on how the participants will manage the research data generated and/or collected during the project
\begin{itemize}
\item What types of data will the project generate/collect?
\item What standards will be used?
\item How will this data be exploited and/or shared/made accessible for verification and re-use? If data cannot be made available, explain why.
\item How will this data be curated and preserved?
\end{itemize}
\end{itemize}

\subsection{Ethics}
If applicable to the nature of work being proposed (working with patients, patient data, or animals),
\begin{itemize}
\item describe how the proposal meets the national legal and ethical requirements of the country or countries where the tasks raising ethical issues are to be carried out; 
\item explain in detail how you intend to address the ethical issues, in particular:
\begin{itemize}
\item research objectives (e.g. study of vulnerable populations, dual use, etc.)
\item research methodology (e.g. clinical trials, involvement of children and related consent procedures, protection of any data collected, etc.) 
\item the potential impact of the research (e.g. dual use issues, environmental damage, stigmatisation of particular social groups, political or financial retaliation, benefit-sharing,  malevolent use, etc.).
\end{itemize}
\end{itemize}

\section{Timeline}
General proposed timeline of both experiments and research dissemination.

\begin{itemize}
\item brief presentation of the overall structure of the work plan;
\item timing of the different work packages and their components (Gantt chart or similar);

\end{itemize}


\section{Budget}
What resources are required to carry out your proposal? Things to consider:
\begin{itemize}
\item equipment and instrumentation
\item costs for third-party to carry out analysis if there is no access to equipment, or a lack of expertise
\item lab space rental
\item personnel costs (who is carrying out the work, how many people)
\item costs for research dissemination: travel to conferences, cost to publish
\end{itemize}


\section{Collaborators}
Are there private companies, research networks, other lab groups that are working on the same topic? Large research projects benefit (and are more likely to be approved) if they are approached in a diverse and multi-disciplinary manner.



\clearpage
%enter info to link bibtex file (.bib) in project
%\bibliography{example}   

%or enter each manually as below: 
\begin{thebibliography}{9}
 
\bibitem{eubel2005}    
Eubel H, Hans-Peter B, Miller AH (2005) Blue-native PAGE in plants: a tool in analysis of protein-protein interactions. Plant Methods 1:11

\end{thebibliography}


\clearpage


\begin{figure}
\centering\includegraphics[width=0.75\linewidth]{example-figure.PNG}
\caption{Example caption.\cite{eubel2005}}
\end{figure}



\end{document} % The document ends here